\begin{conclusion}[Lee--Yang 三次数域的域判别式与分支判别式：坏素数 \texorpdfstring{$37$}{37} 的同步标记]\label{con:xi-terminal-zm-leyang-cubic-field-minimal-discriminant}
沿用结论 \ref{con:xi-terminal-zm-leyang-cubic-field-isomorphism} 的记号，令 \(\beta\) 为 Lee--Yang 三次
\[
P_{\mathrm{LY}}(y)=256y^{3}+411y^{2}+165y+32\in\ZZ[y]
\]
的任一根，并记 \(K:=\QQ(\beta)\)（等价地 \(K=\QQ(\alpha)\)，其中 \(\alpha\) 为 \(16X^{3}-9X^{2}+1=0\) 的根）。
则：
\begin{enumerate}
  \item 多项式判别式与域判别式分别满足
  \[
  \boxed{\ \mathrm{Disc}\bigl(P_{\mathrm{LY}}\bigr)=-3^{9}\cdot 31^{2}\cdot 37\ },
  \qquad
  \boxed{\ \mathrm{Disc}(K)=-999=-3^{3}\cdot 37\ }.
  \]
  并且二者满足指数平方关系
  \[
  \mathrm{Disc}\bigl(P_{\mathrm{LY}}\bigr)=(3^{3}\cdot 31)^2\,\mathrm{Disc}(K).
  \]
  因而素数 \(31\) 仅通过指数 \([\mathcal O_K:\ZZ[\beta]]\) 的平方因子进入 \(\mathrm{Disc}\bigl(P_{\mathrm{LY}}\bigr)\)，并不必作为 \(K/\QQ\) 的分歧素数出现。
  \item 因而扩张 \(K/\QQ\) 的分歧素集恰为 \(\{3,37\}\)；
  其中 \(3\) 为野分歧且全分歧，而 \(37\) 为驯分歧且分解为
  \[
  \boxed{\ 3\mathcal O_K=\mathfrak p_3^{3},\qquad 37\mathcal O_K=\mathfrak p_{37}^{2}\mathfrak q_{37}\ }.
  \]
  特别地，素数 \(37\) 在 \(K\) 中必分歧；结合结论 \ref{con:xi-terminal-zm-elliptic-normalization} 与注记 \ref{rem:xi-terminal-elliptic-37a1-modular-anchor}，其亦为谱椭圆曲线 \(E\) 的导子与 level \(37\) 的模性锚点坏素数，从而在素数层面把 Lee--Yang 分支算术与导出自守对象的最小坏素数严格对齐。
  \item \(K\) 的 \(S_3\)-分裂域之唯一二次分辨子域为
  \[
  \boxed{\ \QQ(\sqrt{-111})\ }.
  \]
  其中 \(-111=-3\cdot 37\) 为 \(\mathrm{Disc}\bigl(P_{\mathrm{LY}}\bigr)\) 的平方自由部分，因此该虚二次域同时携带 Lee--Yang 分支三次与坏素数 \(37\) 的共同判别式指纹。
\end{enumerate}
\end{conclusion}
\begin{proof}[证明要点]
取整生成元 \(\theta:=256\beta\)（从而 \(\QQ(\theta)=\QQ(\beta)\)），则 \(\theta\) 满足首一整式
\(
\theta^{3}+411\theta^{2}+42240\theta+2097152=0
\)。
对该整模型求最大整环并计算域判别式得到 \(\mathrm{Disc}(K)=-999\)，并由判别式判据确定仅在 \(3\) 与 \(37\) 处分歧；
再对 \(p\in\{3,37\}\) 进行素理想分解得到所示分解型。
另一方面，\(\mathrm{Disc}\bigl(P_{\mathrm{LY}}\bigr)\) 可由多项式判别式公式直接计算；并由一般恒等式 \(\mathrm{Disc}\bigl(P_{\mathrm{LY}}\bigr)=\mathrm{Disc}(K)\,[\mathcal O_K:\ZZ[\beta]]^{2}\) 得到上式指数平方关系。
可复现核验脚本：\texttt{scripts/exp\_fold\_zm\_elliptic\_leyang\_arithmetic\_audit.py}。证毕。
\end{proof}

\begin{conjecture}[射线类场刻画：\texorpdfstring{$\QQ(\sqrt{-111})$}{Q(sqrt(-111))} 上的循环三次分裂域]\label{conj:xi-terminal-zm-leyang-ray-class-field}
令 \(L\) 为 Lee--Yang 三次 \(P_{\mathrm{LY}}(y)\) 的分裂域，并令
\(
K_{\Delta}:=\QQ(\sqrt{-111})
\)
为其二次分辨子域（等价地，\(K_{\Delta}=\QQ(\sqrt{\mathrm{Disc}(P_{\mathrm{LY}})})\)，见结论 \ref{con:xi-terminal-zm-leyang-cubic-arithmetic}）。
则 \(L/K_{\Delta}\) 为循环三次扩张。
进一步猜想 \(L\) 恰为 \(K_{\Delta}\) 的某一射线类域，并且其导子可取为
\[
\mathfrak f=3^{4}\cdot 31
\qquad(\text{按 }K_{\Delta}\text{ 的理想分解理解}).
\]
等价地，记 \(y_1,y_2,y_3\) 为重量覆盖 \(\pi_y:E\to\mathbb P^1_y\) 的三条非平凡有限分歧值（即 \(P_{\mathrm{LY}}(y)=0\) 的三根），则
\[
L=\QQ(y_1,y_2,y_3)=K_{\Delta}(y_1,y_2,y_3).
\]
\end{conjecture}

\begin{remark}[判别式分解与导子候选]\label{rem:xi-terminal-zm-leyang-rayclass-discriminant-motivation}
由定理 \ref{thm:xi-terminal-zm-leyang-elliptic-structure} 的判别式分解
\[
\mathrm{Disc}\bigl(P_{\mathrm{LY}}(y)\bigr)=-3^{9}\cdot 31^{2}\cdot 37
=-111\cdot(3^{4}\cdot 31)^{2}
\]
可见二次分辨子域的基本判别式 \(-111\) 被平方因子精确剥离；在类场论的导子--判别式对照中，这一平方结构为 \(\mathfrak f\) 的候选提供了自然的可审计接口。
\end{remark}

\begin{conclusion}[分歧特征值与分歧参数的二次互换：缩放生成元的同域性与指数公式]\label{con:xi-terminal-zm-leyang-uv-quadratic-exchange-index}
仍取非平凡分歧值 \(y_0\)（即 \(P_{\mathrm{LY}}(y_0)=0\)）与其对应的并合谱值 \(\lambda_0\)（即 \(\Pi(\lambda_0,y_0)=\partial_\lambda\Pi(\lambda_0,y_0)=0\)，等价于 \(16\lambda_0^3-9\lambda_0^2+1=0\) 且 \(y_0=\frac{4\lambda_0^3-3\lambda_0^2-2\lambda_0+1}{4\lambda_0}\)，见结论 \ref{con:xi-terminal-zm-leyang-ramification-critical-values}）。
定义缩放生成元
\[
u:=16\lambda_0,\qquad v:=256y_0.
\]
则 \(u\) 与 \(v\) 生成同一三次域 \(K\)，并且满足显式二次互换关系
\begin{equation}\label{eq:xi-terminal-zm-leyang-uv-quadratic-exchange}
\boxed{\;
3u^2+(v+128)u-768=0\;}
\end{equation}
从而 \(u\in\QQ(v)\) 且 \(v\in\QQ(u)\)，并且 \(\QQ(u)=\QQ(v)=K\)。
等价地，利用 \(u\neq 0\) 可写出
\[
v=u^2-12u-128+\frac{1024}{u}=-3u^2+24u-128\in\ZZ[u],
\]
因而 \(\ZZ[v]\subseteq \ZZ[u]\)。
\par
记
\[
f_u(t):=t^3-9t^2+256,\qquad
f_v(t):=t^3+411t^2+42240t+2^{21}.
\]
则 \(u\) 为 \(f_u\) 的根且 \(v\) 为 \(f_v\) 的根；并且二者的判别式满足
\[
\mathrm{Disc}(f_u)=-2^{10}\cdot 3^{3}\cdot 37,\qquad
\mathrm{Disc}(f_v)=-2^{16}\cdot 3^{9}\cdot 31^{2}\cdot 37.
\]
从而两单生成整环的指数满足严格公式
\begin{equation}\label{eq:xi-terminal-zm-leyang-uv-index}
\boxed{\;
[\ZZ[u]:\ZZ[v]]=2^3\cdot 3^3\cdot 31=6696.\;}
\end{equation}
\end{conclusion}
\begin{proof}[证明要点]
由结论 \ref{con:xi-terminal-zm-leyang-ramification-critical-values} 的参数化
\(
y_0=\frac{4\lambda_0^3-3\lambda_0^2-2\lambda_0+1}{4\lambda_0}
\)
得
\[
v
=256y_0
=\frac{64}{\lambda_0}\bigl(4\lambda_0^3-3\lambda_0^2-2\lambda_0+1\bigr)
=u^2-12u-128+\frac{1024}{u}.
\]
又由 \(16\lambda_0^3-9\lambda_0^2+1=0\) 得 \(u\) 满足 \(u^3-9u^2+256=0\)，从而 \(\frac{1024}{u}=-4u^2+36u\)，得到 \(v=-3u^2+24u-128\in\ZZ[u]\) 及包含 \(\ZZ[v]\subseteq\ZZ[u]\)。
将 \(u^3=9u^2-256\) 代入恒等式 \(vu=u^3-12u^2-128u+1024\) 并整理即得 \eqref{eq:xi-terminal-zm-leyang-uv-quadratic-exchange}，从而 \(\QQ(u)=\QQ(v)\)。
\par
由 \(\lambda_0=u/16\) 代入三次方程 \(16\lambda^3-9\lambda^2+1=0\) 得 \(f_u(u)=0\)；
由 \(y_0=v/256\) 代入 \(P_{\mathrm{LY}}(y)=0\) 并清分母得 \(f_v(v)=0\)。
判别式由缩放公式与 \(\mathrm{Disc}(16\lambda^3-9\lambda^2+1)\)、\(\mathrm{Disc}(P_{\mathrm{LY}})\) 的闭式（定理 \ref{thm:xi-terminal-zm-leyang-elliptic-structure}）直接推出。
最后，由 \(\mathrm{Disc}(\ZZ[\theta])=\mathrm{Disc}(m_\theta)\)（\(\theta\) 为代数整数且 \(m_\theta\) 为其首一极小多项式）与
\(
\mathrm{Disc}(\ZZ[v])=\mathrm{Disc}(\ZZ[u])\,[\ZZ[u]:\ZZ[v]]^{2}
\)
得
\(
[\ZZ[u]:\ZZ[v]]=\sqrt{\mathrm{Disc}(f_v)/\mathrm{Disc}(f_u)}=2^3\cdot 3^3\cdot 31
\)。
证毕。
\end{proof}

\begin{remark}[二次分辨子域的规范特征与权一自守对接]\label{rem:xi-terminal-zm-leyang-artin-weight1-interface}
令 \(L\) 为 \(P_{\mathrm{LY}}\) 的分裂域。
由 \(\mathrm{Disc}\bigl(P_{\mathrm{LY}}\bigr)=-3^{9}\cdot 31^{2}\cdot 37\) 的非平方性可知 \(\mathrm{Gal}(L/\QQ)\cong S_3\)，并且其唯一二次分辨子域为
\(
\QQ(\sqrt{-111})
\)
（结论 \ref{con:xi-terminal-zm-leyang-cubic-field-minimal-discriminant}）。
取 \(S_3\hookrightarrow \mathrm{GL}_2(\CC)\) 为标准二维不可约表示，则复合
\[
\rho_{\mathrm{LY}}:\mathrm{Gal}(\overline{\QQ}/\QQ)\twoheadrightarrow \mathrm{Gal}(L/\QQ)\hookrightarrow \mathrm{GL}_2(\CC)
\]
给出一个二维不可约 Artin 表示，其行列式等于符号特征，从而与二次特征 \(\chi_{-111}\) 对齐。
该表示的分歧素集包含于 \(\{3,31,37\}\)，其中 \(31\) 由 \(\mathrm{Disc}\bigl(P_{\mathrm{LY}}\bigr)\) 的指数平方因子进入；而在域判别式口径下仅由 \(\{3,37\}\) 触发（\(\mathrm{Disc}(K)=-999\)）。
由于 \(S_3\) 为可解群像，\(\rho_{\mathrm{LY}}\) 属于经典 dihedral/可解像 reciprocity 的已知范畴；
从而其 Artin \(L\)-函数可与权 \(1\) 模形式的 \(L\)-函数对接，并据此将 Lee--Yang 分歧三次转写为一条可计算的权一自守数据通道。\cite{DeligneSerre1974WeightOne}
\end{remark}

\begin{remark}[\texorpdfstring{$\sqrt{-3}$}{sqrt(-3)} 自反通道与 \texorpdfstring{$\sqrt{37}$}{sqrt(37)} 模性锚点的二次扭转：\texorpdfstring{$\QQ(\sqrt{-111})$}{Q(sqrt(-111))} 的出现]\label{rem:xi-terminal-zm-leyang-sqrtminus111-twist}
结论 \ref{con:xi-terminal-zm-self-inversive-unique} 表明自反性仅在 \(y^2+y+1=0\) 时发生，从而唯一非平凡自反通道由虚二次域
\(
K_{-3}:=\QQ(\sqrt{-3})
\)
标记。
另一方面，结论 \ref{con:xi-terminal-37-squarefree-coupling} 给出椭圆锚点的 \(2\)-division 伴随二次域为
\(
K_{37}:=\QQ(\sqrt{37})
\)。
在双二次扩张 \(K_{37,-3}:=\QQ(\sqrt{37},\sqrt{-3})\) 中，第三个二次子域为
\[
\QQ(\sqrt{37}\sqrt{-3})=\QQ(\sqrt{-111}),
\]
其与 Lee--Yang 三次的二次分辨子域一致（结论 \ref{con:xi-terminal-zm-leyang-cubic-field-minimal-discriminant}）。
等价地，对任意 \(p\nmid 3\cdot 37\) 有 Kronecker 符号的乘法恒等式
\[
\left(\frac{-111}{p}\right)=\left(\frac{37}{p}\right)\left(\frac{-3}{p}\right),
\]
从而 Lee--Yang 三次分歧的二次特征可视为在 level \(37\) 椭圆锚点之上施加的唯一 \((-3)\) 扭转读出。
\end{remark}

